\documentclass[aspectratio=169]{beamer}

\usepackage{fontspec}
\usepackage{amsmath}
\usepackage{booktabs}
\usepackage{listings}
\usepackage{xcolor}

\usetheme{Madrid}
\setmainfont{DejaVu Sans}
\setsansfont{DejaVu Sans}
\setmonofont{DejaVu Sans Mono}

\definecolor{codebg}{RGB}{245,245,245}
\lstset{
  basicstyle=\ttfamily\scriptsize,
  backgroundcolor=\color{codebg},
  breaklines=true,
  frame=single,
  columns=fullflexible
}

\title{Cơ chế thay đổi kích thước ảnh PKA}
\subtitle{\texttt{prepare\_all\_pka.py} và \texttt{prepare\_pka\_sample.py}}
\author{P2VG preprocessing}
\date{\today}

\begin{document}

\begin{frame}
  \titlepage
\end{frame}

\begin{frame}{Mục tiêu tiền xử lý}
  \begin{itemize}
    \item Đưa các ảnh MRI PKA về định dạng thống nhất trước khi train model.
    \item Chuẩn hóa spacing voxel về \((1.0, 1.0, 4.0)\) mm.
    \item Chuẩn hóa kích thước mảng ảnh về \(256 \times 256 \times 32\).
    \item Xử lý cả hai modality mặc định: \texttt{fused} và \texttt{axt2}.
  \end{itemize}
\end{frame}

\begin{frame}{Vai trò của \texttt{prepare\_all\_pka.py}}
  \begin{itemize}
    \item File này không trực tiếp resize ảnh.
    \item Nó tìm tất cả subject có file \texttt{sub-*\_fused.nii.gz}.
    \item Với mỗi subject, nó gọi \texttt{prepare\_pka\_sample.py}.
    \item Xử lý song song bằng \texttt{multiprocessing.Pool(8)}.
  \end{itemize}

  \vspace{0.4cm}
  \begin{lstlisting}
--mode preserve_axis
--xy_size 256
--depth_size 32
  \end{lstlisting}
\end{frame}

\begin{frame}{Luồng xử lý tổng quát}
  \begin{enumerate}
    \item Đọc ảnh NIfTI bằng \texttt{nibabel}.
    \item Lấy dữ liệu ảnh dạng \texttt{float32}.
    \item Resample theo spacing mục tiêu.
    \item Resize trục sâu về 32 slices.
    \item Resize hai trục mặt phẳng ảnh về \(256 \times 256\).
    \item Cập nhật affine để phản ánh spacing mới.
    \item Lưu ảnh mới và file metadata JSON.
  \end{enumerate}
\end{frame}

\begin{frame}{Bước 1: Resample theo spacing vật lý}
  \texttt{prepare\_pka\_sample.py} lấy spacing cũ từ header:

  \[
    old\_zooms = (s_x, s_y, s_z)
  \]

  Spacing mục tiêu mặc định:

  \[
    target\_zooms = (1.0, 1.0, 4.0)
  \]

  Hệ số zoom được tính:

  \[
    zoom\_factors = \frac{old\_zooms}{target\_zooms}
  \]

  Sau đó nội suy tuyến tính:

  \[
    data' = zoom(data, zoom\_factors, order=1)
  \]
\end{frame}

\begin{frame}{Ý nghĩa của resample spacing}
  \begin{itemize}
    \item Đây là bước chuẩn hóa theo kích thước voxel thật trong không gian vật lý.
    \item Nếu voxel ban đầu dày/mỏng khác nhau, số voxel sau resample sẽ thay đổi.
    \item Mục tiêu là để các ảnh có cùng độ phân giải vật lý xấp xỉ nhau.
    \item \texttt{order=1} nghĩa là dùng nội suy tuyến tính, phù hợp cho ảnh MRI intensity.
  \end{itemize}
\end{frame}

\begin{frame}{Bước 2: Chuẩn hóa depth về 32}
  Sau khi resample spacing, script kiểm tra trục cuối:

  \[
    D = final\_data.shape[2]
  \]

  Nếu \(D \ne 32\), hệ số zoom theo depth là:

  \[
    depth\_zoom = \frac{32}{D}
  \]

  Và resize:

  \[
    final\_data = zoom(final\_data, (1.0, 1.0, depth\_zoom), order=1)
  \]
\end{frame}

\begin{frame}{Bước 3: Chuẩn hóa mặt phẳng XY về 256}
  Nếu chiều \(X\) hoặc \(Y\) chưa bằng 256:

  \[
    xy\_zoom_x = \frac{256}{X}
  \]

  \[
    xy\_zoom_y = \frac{256}{Y}
  \]

  Sau đó resize:

  \[
    final\_data = zoom(final\_data, (xy\_zoom_x, xy\_zoom_y, 1.0), order=1)
  \]

  Trục depth lúc này được giữ nguyên.
\end{frame}

\begin{frame}{Vì sao cần cập nhật affine?}
  \begin{itemize}
    \item Khi thay đổi số voxel, spacing vật lý tương ứng cũng thay đổi.
    \item Script tính lại:
  \end{itemize}

  \[
    final\_zooms = target\_zooms \times
    \frac{shape\_before\_normalize}{shape\_after\_normalize}
  \]

  \begin{itemize}
    \item Sau đó affine được rescale bằng hàm \texttt{\_rescaled\_affine}.
    \item Nhờ vậy file NIfTI đầu ra vẫn có thông tin không gian hợp lý.
  \end{itemize}
\end{frame}

\begin{frame}{Hai chế độ xử lý}
  \begin{table}
    \centering
    \begin{tabular}{lll}
      \toprule
      Chế độ & Cách xử lý & Khi dùng \\
      \midrule
      \texttt{canonical} & Đưa về orientation RAS rồi resample spacing & Chuẩn hóa hướng ảnh \\
      \texttt{preserve\_axis} & Giữ thứ tự axis gốc, resample và resize & Đang dùng trong pipeline hiện tại \\
      \bottomrule
    \end{tabular}
  \end{table}

  \vspace{0.3cm}
  \texttt{prepare\_all\_pka.py} gọi chế độ \texttt{preserve\_axis}.
\end{frame}

\begin{frame}{Kết quả đầu ra}
  Với mỗi subject và modality, script lưu:

  \begin{itemize}
    \item File ảnh mới: \texttt{sub-XXXX\_fused.nii.gz}, \texttt{sub-XXXX\_axt2.nii.gz}.
    \item Metadata JSON: mô tả shape, spacing, orientation trước và sau xử lý.
  \end{itemize}

  \vspace{0.3cm}
  Kích thước mảng mục tiêu trong pipeline:

  \[
    256 \times 256 \times 32
  \]
\end{frame}

\begin{frame}{Tóm tắt ngắn gọn}
  \begin{itemize}
    \item \texttt{prepare\_all\_pka.py}: chạy batch cho toàn bộ subject.
    \item \texttt{prepare\_pka\_sample.py}: thực hiện resize/resample thật sự.
    \item Đầu tiên chuẩn hóa spacing vật lý bằng nội suy tuyến tính.
    \item Sau đó ép depth về 32 và XY về \(256 \times 256\).
    \item Cuối cùng cập nhật affine và lưu metadata để kiểm tra lại.
  \end{itemize}
\end{frame}

\end{document}
